\chapter{The Final Conversation: Resume Drill, Company Knowledge and HR}
\companion{InfoEdge DS interview experience: HR round and ``Why InfoEdge'' (at 12:26)}{https://youtu.be/nWyBQhGnHaM?t=746}

\section{Everything on your resume will be drilled}
The hire's interviewers followed his projects (PCA, t-SNE) into depth. For every project and technique on your resume prepare:
\begin{enumerate}
\item A 60-second story: problem, why it mattered, data, approach, your specific contribution, result with a number, what you would do differently.
\item The maths of every method you name, at the level of this book.
\item The baseline you beat and how you evaluated (split, metric, significance).
\item One failure or surprise and what you learned.
\end{enumerate}

\section{Company knowledge}
InfoEdge (India) Ltd: founded in the mid-1990s by Sanjeev Bikhchandani; Naukri.com (launched 1997, India's largest job portal), Jeevansathi.com (matrimony), 99acres.com
(real estate), Shiksha.com (education), and an investment arm with stakes in companies such as Zomato and PolicyBazaar. Listed on Indian stock exchanges in 2006. \emph{Verify
years and figures from InfoEdge's annual report before quoting them.} From the 2026--27 talk: 100 crore+ applies a year through the job-posting engine; Resdex serves about 60K
recruiters and 16 lakh searches a day; the Data Factory runs 400+ AI models; GenAI work includes in-house fine-tuned and distilled LLMs, voice agents and the AI REX agentic hiring
product. Work model (as reported): mostly in-office.

\section{``Why InfoEdge?''}
The hire called this the most important HR question. A strong answer has three specific parts:
\begin{enumerate}
\item \textbf{Real problems at scale}: ML that touches crores of job seekers and lakhs of recruiters; search, ranking, matching and recommendation are core, not side projects.
\item \textbf{Technical depth}: cite something from the pre-placement talk (hybrid retrieval in PremiumX, the SilverSkAI taxonomy, AI REX's agents, in-house LLMs) and why it excites you.
\item \textbf{Breadth and business}: several products (jobs, real estate, matrimony, education) let you learn many problem types; the company wants data scientists who understand
  business, which matches what you want to become.
\end{enumerate}

\section{Questions to ask them}
How do teams decide between a classical ranker and an LLM-based approach given latency and cost? How is model impact measured (A/B culture)? How do new data scientists ramp up?
What does a successful first six months look like?

\stuck{InfoEdge DS interview experience (full video)}{\ytid{nWyBQhGnHaM}}
\section{Questions}

\begin{qa}{Tell me about yourself (for a data scientist role at InfoEdge).}
\oneline{A 60--90 second story: who you are academically, the two or three ML experiences most relevant to search, ranking, NLP or tabular modelling with one concrete result each, and why that
leads you to InfoEdge now.}
\why Example structure: ``I am an M.Tech AI student at IIT Jodhpur. In my project on X I built Y, which improved Z by N\%. I enjoy problems where models meet users at scale, which is why the
matching and search work you presented, especially hybrid retrieval in PremiumX, is exactly what I want to work on.''
\end{qa}

\begin{qa}{Describe a project where your model did not work as expected. What did you do?}
\oneline{Pick a real case, explain the symptom (for example great offline, poor online), how you diagnosed it (leakage, distribution shift, wrong metric), what you changed, and what you now do by
default because of it.}
\end{qa}

\begin{qa}{How do you keep up with ML, and what recent development do you find most useful for a company like InfoEdge?}
\oneline{Name concrete sources and one development with a clear application: for example retrieval-augmented generation with hybrid search for grounded job-discovery chat, or distillation of
LLMs into small models to serve ranking features cheaply.}
\end{qa}
